\begin{proof}
This is a scalar contraction-completion result; we provide its elementary
Schur proof rather than relying on an external completion theorem (compare
the classical mechanism discussed in~\cite{parrott2025}). Write
\[
 H=\begin{pmatrix}h&b\\c&D\end{pmatrix}
\]
relative to $\operatorname{span}\{p\}\oplus p^\perp$. Since
$\Fnorm D<\gamma$, the complement block
\[
 \mathcal L=\begin{pmatrix}\gamma I&D\\D^*&\gamma I\end{pmatrix}
\]
is strictly positive. Put $t=\gamma^2$ and $x=h-\tau$.
The norm inequality for $M=\left(\begin{smallmatrix}x&b\\c&D\end{smallmatrix}\right)$
is equivalent to positivity of
$\left(\begin{smallmatrix}\gamma I&M\\M^*&\gamma I\end{smallmatrix}\right)$.
Its Schur complement with respect to $\mathcal L$ is the scalar block
\[
 \begin{pmatrix}
 \gamma\{1-b(tI-D^*D)^{-1}b^*\}
 &x+bD^*(tI-DD^*)^{-1}c\\
 \overline x+c^*(tI-DD^*)^{-1}Db^*
 &\gamma\{1-c^*(tI-DD^*)^{-1}c\}
 \end{pmatrix}.
\]
The given column and row norm bounds make both diagonal entries
nonnegative. Choose $x$ to make the off-diagonal zero. This proves
existence. If the column norm equals $\gamma$, its scalar Schur
complement, the first diagonal entry, is zero; positivity then forces the
off-diagonal to vanish, proving uniqueness. All formulas preserve real
data.
\end{proof}
